In Sec~\ref{sec:exp-verify}, we first empirically verify the theoretical analysis relating to the Wasserstein upper bounds. We then evaluate the performance of different sampling algorithms on standard image generation benchmarks, including CIFAR-10~\cite{Krizhevsky2009LearningML} and ImageNet $64\times 64$~\cite{Deng2009ImageNetAL} in Sec~\ref{sec:exp-standard}. Lastly, we employ Restart on text-to-image generation, using Stable Diffusion model~\cite{Rombach2021HighResolutionIS} pre-trained on LAION-5B~\cite{Schuhmann2022LAION5BAO} with resolution $512\times 512$, in Sec~\ref{exp-t2i}.

\subsection{Additional sampling error versus contracted error}
\label{sec:exp-verify}
Our proposed Restart sampling algorithm demonstrates a higher contraction effect and smaller addition sampling error compared to SDE, according to Theorem~\ref{thm1_informal} and Theorem~\ref{thm2_informal}. Although our theoretical analysis compares the upper bounds of the total, contracted and additional sampling errors, we further verify their relative values through a synthetic experiment.

\textbf{Setup } We construct a $20$-dimensional dataset with 2000 points sampled from a Gaussian mixture, and train a four-layer MLP to approximate the score field $\nabla_x \log p_t$. We implement the ODE, SDE, and Restart methods within a predefined time range of $[\tmin, \tmax]=[1.0, 1.5]$, where the process outside this range is conducted via the first-order ODE. To compute various error types, we define the distributions generated by three methods as outlined in the proof of Theorem~\ref{thm2_informal} and directly gauge the errors at end of simulation $t=0$ instead of $t=\tmin$: (1) the generated distribution as {\footnotesize $p^{\textrm{Sampler}}_{0}$}, where $\textrm{Sampler} \in {\footnotesize \{\bODE, \bSDE, \bRestart(K)\}}$; (2) an auxiliary distribution {\footnotesize $q^{\textrm{Sampler}}_{0}$} initiating from true distribution $p_{\tmax}$ at time $\tmax$. The only difference between {\footnotesize $p^{\textrm{Sampler}}_{0}$} and {\footnotesize $q^{\textrm{Sampler}}_{0}$} is their initial distribution at $\tmax$~({\footnotesize $p^{\bODE}_{\tmax}$} versus $p_{\tmax}$); and (3) the true data distribution $p_{0}$. In line with Theorem~\ref{thm2_informal}, we use Wasserstein-1 distance {\footnotesize $W_1(p^{\textrm{Sampler}}_{0},q^{\textrm{Sampler}}_{0})$ / $W_1(q^{\textrm{Sampler}}_{0}, p_{0})$} to measure the contracted error / additional sampling error, respectively. Ultimately, the total error corresponds to {\footnotesize $W_1(p^{\textrm{Sampler}}_{0}, p_{0})$}. Detailed information about dataset, metric and model can be found in the Appendix~\ref{app:syn}.


\textbf{Results } In our experiment, we adjust the parameters for all three processes and calculate the total, contracted, and additional sampling errors across all parameter settings. Figure \ref{fig:approx_contract} depicts the Pareto frontier of additional sampling error versus contracted error. We can see that Restart consistently achieves lower contracted error for a given level of additional sampling error, compared to both the ODE and SDE methods, as predicted by theory. In Figure \ref{fig:approx_fid}, we observe that the Restart method obtains a smaller total error within the additional sampling error range of $[0.8, 0.85]$. During this range, Restart also displays a strictly reduced contracted error, as illustrated in Figure \ref{fig:approx_contract}. This aligns with our theoretical analysis, suggesting that the Restart method offers a smaller total error due to its enhanced contraction effects. From Figure \ref{fig:nfe_total}, Restart also strikes an better balance between efficiency and quality, as it achieves a lower total error at a given NFE.

\begin{figure*}[t]
    \centering
    \subfigure[]{\includegraphics[width=0.33\textwidth]{img/approx_vs_contracted_no_sde.pdf}\label{fig:approx_contract}}\hfill
   \subfigure[]{\includegraphics[width=0.325\textwidth]{img/approx_vs_fid_scatter_no_sde.pdf}\label{fig:approx_fid}}\hfill
  \subfigure[]{\includegraphics[width=0.32\textwidth]{img/NFE_vs_FID_no_sde.pdf}\label{fig:nfe_total}}
  \vspace{-6pt}
    \caption{Additional sampling error versus \textbf{(a)} contracted error, where the Pareto frontier is plotted and \textbf{(b)} total error, where the scatter plot is provided. \textbf{(c)} Pareto frontier of NFE versus total error.}
    \vspace{-12pt}
\end{figure*}


\subsection{Experiments on standard benchmarks}
\label{sec:exp-standard}
\begin{figure*}[h]
\vspace{-6pt}
  \centering
  \subfigure[CIFAR-10, VP]{\includegraphics[width=0.45\textwidth]{img/cifarvp_fid50k.pdf}\label{fig:cifar10}}\hfill
  \subfigure[ImageNet $64{\times} 64$, EDM]{\includegraphics[ width=0.505\textwidth, trim = 0.9cm 0cm 0.9cm 0cm]{img/imagenet_fid50k.pdf}\label{fig:imagenet}}
  \vspace{-3pt}
  \caption{FID versus NFE on \textbf{(a)} unconditional generation on CIFAR-10 with VP; \textbf{(b)} class-conditional generation on ImageNet with EDM.}
  \label{fig:example}
  \vspace{-5pt}
\end{figure*}

\begin{figure*}
\vspace{-3pt}
\centering
    \includegraphics[width=0.5\textwidth]{img/cifarvp_dpm.pdf}
    \caption{CIFAR-10, VP, in the low NFE regime. Restart consistently outperforms the DPM-solver with an NFE ranging from 16 to 36.}
        \label{fig:cifar-dpm}
    \vspace{-15pt}
\end{figure*}
% \begin{figure*}[t]
%   \centering
%   \subfigure[]{\includegraphics[width=0.3\textwidth]{img/cifar_k.pdf}\label{fig:k}}\hfill
% \subfigure[]{\includegraphics[width=0.3\textwidth]{img/cifar_k.pdf}\label{fig:k}}
%   \caption{}
%   \label{fig:example}
% \end{figure*}

% \begin{minipage}{\textwidth}
%   \begin{minipage}[b]{0.65\textwidth}
%     \centering
%     % \rule{6.4cm}{3.6cm}
%     \includegraphics[width=0.5\textwidth]{img/cifar_k.pdf}\label{fig:k}
%     \captionof{figure}{A table beside a figure}
%   \end{minipage}
%   \hfill
%   \begin{minipage}[b]{0.3\textwidth}
%     \centering
%     \scriptsize
%     \captionof{table}{Unconditional generation on CIFAR-10 using EDM~\cite{Karras2022ElucidatingTD}}
%     \begin{tabular}{l c c c}
%     \toprule
%          & NFE & FID  \\
%          \midrule \textit{EDM-VP}~\cite{Karras2022ElucidatingTD}\\ \midrule
% {ODE~(Heun)}& 63 & 1.97\\ & 35& 1.97 \\
%          {Improved SDE}~\cite{Karras2022ElucidatingTD}
%          & 63& 2.27 \\
%          & 35& 2.45 \\
%         {Restart} & 43& \textbf{1.90}\\
%         \midrule \textit{PFGM++}~\cite{Xu2023PFGMUT}\\ \midrule
%         {ODE~(Heun)}& 63 & 1.91\\ 
%         & 35& 1.91 \\
%         Restart & 43 & \textbf{1.88}\\
%         \bottomrule
%     \end{tabular}
%       %\captionof{table}{A table beside a figure}
%     \end{minipage}
%        \label{tab:cifar10-edm}
% \end{minipage}

To evaluate the sample quality and inference speed, we report the FID score~\cite{Heusel2017GANsTB}~(lower is better) on 50K samplers and the number of function evaluations~(NFE). We borrow the pretrained VP/EDM/PFGM++ models on CIFAR-10 or ImageNet $64\times 64$ from \cite{Karras2022ElucidatingTD, Xu2023PFGMUT}. We also use the EDM discretization scheme~\cite{Karras2022ElucidatingTD}~(see Appendix~\ref{app:edm-time} for details) during sampling. 




For the proposed Restart sampler, the hyperparameters include the number of steps in the main/Restart backward processes, the number of Restart iteration $K$, as well as the time interval $[\tmin, \tmax]$. 
We pick the $\tmin$ and $\tmax$ from the list of time steps in EDM discretization scheme with a number of steps $18$. For example, for CIFAR-10~(VP) with NFE${=}75$, we choose $\tmin{=}0.06,\tmax{=}0.30,  K{=}10$, where $0.30$/$0.06$ is the $12^{\textrm{th}}$/$14^{\textrm{th}}$ time step in the EDM scheme. We also adopt EDM scheme for the Restart backward process in $[\tmin, \tmax]$. In addition, we apply the multi-level Restart strategy~(Sec~\ref{sec:practical}) to mitigate the error at early time steps for the more challenging ImageNet $64\times 64$. We provide the detailed Restart configurations in Appendix~\ref{app:config}. 

% \begin{wrapfigure}{r}{0.35\textwidth} 
% \vspace{-11pt}
% \centering
% \scriptsize
% \captionof{table}{Unconditional generation on CIFAR-10 using EDM~\cite{Karras2022ElucidatingTD}}
%     \begin{tabular}{l c c c}
%     \toprule
%          & NFE & FID  \\
%          \midrule \textit{EDM-VP}~\cite{Karras2022ElucidatingTD}\\ \midrule
% {ODE~(Heun)}& 63 & 1.97\\ & 35& 1.97 \\
%          {Improved SDE}~\cite{Karras2022ElucidatingTD}
%          & 63& 2.27 \\
%          & 35& 2.45 \\
%         {Restart} & 43& \textbf{1.90}\\
%         \midrule \textit{PFGM++}~\cite{Xu2023PFGMUT}\\ \midrule
%         {ODE~(Heun)}& 63 & 1.91\\ 
%         & 35& 1.91 \\
%         Restart & 43 & \textbf{1.88}\\
%         \bottomrule
%     \end{tabular}
%     \label{tab:cifar10-edm}
%     \vspace{-5pt}
% \end{wrapfigure}

For SDE, we compare with the previously best-performing stochastic samplers proposed by \cite{Karras2022ElucidatingTD} (\textbf{Improved SDE}). We use their optimal hyperparameters for each dataset. We also report the FID scores of the adaptive SDE~\cite{JolicoeurMartineau2021GottaGF}~(\textbf{Gonna Go Fast}) on CIFAR-10~(VP). Since the vanilla reverse-diffusion SDE~\cite{Song2020ScoreBasedGM} has a significantly higher FID score, we omit its results from the main charts and defer them to Appendix~\ref{app:extra-exp}. For ODE samplers, we compare with the Heun's $2^{\textrm{nd}}$ order method~\cite{Ascher1998ComputerMF} (\textbf{Heun}), which arguably provides an excellent trade-off between discretization errors and NFE~\cite{Karras2022ElucidatingTD}. To ensure a fair comparison, we use Heun's method as the sampler in the main/Restart backward processes in Restart. 


We report the FID score versus NFE in Figure~\ref{fig:cifar10} and Table~\ref{tab:cifar10-edm} on CIFAR-10, and Figure~\ref{fig:imagenet} on ImageNet $64\times 64$ with EDM. Our main findings are: \textbf{(1)} Restart outperforms other SDE or ODE samplers in balancing quality and speed, across datasets and models. As demonstrated in the figures, Restart achieves a 10-fold / 2-fold acceleration compared to previous best SDE results on CIFAR-10~(VP) / ImageNet $64\times 64$~(EDM) at the same FID score. In comparison to ODE sampler~(Heun), Restart obtains a better FID score, with the gap increasing significantly with NFE. \textbf{(2)} For stronger models such as EDM and PFGM++, Restart further improve over the ODE baseline on CIFAR-10. In contrast, the Improved SDE negatively impacts performance of EDM, as also observed in \cite{Karras2022ElucidatingTD}. It suggests that Restart incorporates stochasticity more effectively. \textbf{(3)} Restart establishes new state-of-the-art FID scores for UNet architectures without additional training. In particular, Restart achieves FID scores of 1.36 on class-cond. ImageNet $64\times 64$ with EDM, and 1.88 on uncond. CIFAR-10 with PFGM++.


\textbf{To further validate that Restart can be applied in low NFE regime}, we show that one can employ faster ODE solvers such as the \textbf{DPM-solver-3}~\cite{lu2022dpm} to further accelerate Restart. \Figref{fig:cifar-dpm} shows that the Restart consistently outperforms the DPM-solver with an NFE ranging from 16 to 36. This demonstrates Restart's capability to excel over ODE samplers, even in the small NFE regime. It also suggests that Restart can consistently improve other ODE samplers, not limited to the DDIM, Heun. Surprisingly, when paired with the DPM-solver, \textbf{Restart achieves an FID score of $2.11$ on VP setting when NFE is 30, which is significantly lower than any previous numbers (even lower than the SDE sampler with an NFE greater than 1000
 in \cite{Song2020ScoreBasedGM}), and make VP model on par with the performance with more advanced models (such as EDM).} We include detailed Restart configuration in Table~\ref{tab:app-cifar-restart} in Appendix~\ref{app:config}.

 \begin{wrapfigure}{r}{0.52\textwidth} 
\vspace{-10pt}
	\begin{minipage}{0.4\linewidth}
		\captionof{table}{\footnotesize Uncond. CIFAR-10 with EDM and PFGM++}
		\label{tab:cifar10-edm}
		\centering
  \scriptsize
    \begin{tabular}{l @{\hspace{0.5\tabcolsep}} c @{\hspace{0.5\tabcolsep}} c}
    \toprule
         & NFE & FID  \\
         \midrule \textit{EDM-VP}~\cite{Karras2022ElucidatingTD}\\ \midrule
{ODE~(Heun)}& 63 & 1.97\\ & 35& 1.97 \\
         {Improved SDE}
         & 63& 2.27 \\
         & 35& 2.45 \\
        {Restart} & 43& \textbf{1.90}\\
        \midrule \textit{PFGM++}~\cite{Xu2023PFGMUT}\\ \midrule
        {ODE~(Heun)}& 63 & 1.91\\ 
        & 35& 1.91 \\
        Restart & 43 & \textbf{1.88}\\
        \bottomrule
    \end{tabular}
	\end{minipage}\hfill
	\begin{minipage}{0.55\linewidth}
    \includegraphics[width=0.95\textwidth]{img/cifar_k.pdf}
    %\subfigure[]{\includegraphics[width=0.45\textwidth]{img/cifar_t_min.pdf}\label{fig:tmin}}
    \vspace{-3pt}
    \captionof{figure}{\footnotesize FID score with a varying number of Restart iterations $K$.}
    \label{fig:k}
	\end{minipage}
    \vspace{-10pt}
\end{wrapfigure}


Theorem~\ref{thm2} shows that each Restart iteration reduces the contracted errors while increasing the additional sampling errors in the backward process. In \Figref{fig:k}, we explore the choice of the number of Restart iterations $K$ on CIFAR-10. We find that FID score initially improves and later worsens with increasing iterations $K$, with a smaller turning point for stronger EDM model. This supports the theoretical analysis that sampling errors will eventually outweigh the contraction benefits as $K$ increases, and EDM only permits fewer Restart iterations due to smaller accumulated errors. It also suggests that, as a rule of thumb, we should apply greater Restart strength (\eg larger $K$) for weaker or smaller architectures and vice versa. 

% We also investigate the impact of varying $\tmin$ when $\tmax=\tmin+0.3$. \Figref{fig:tmin} reveals that FID scores achieve a minimum at a $\tmin$ close to $0$ on VP, indicating higher accumulated errors at the end of sampling and poor neural estimations at small $t$. Note that the Restart interval $0.3$ is about twice the length of the one in Table~\ref{tab:cifar10-edm} and Restart does not outperform the ODE baseline on EDM. This suggests that, as a rule of thumb, we should apply greater Restart strength (\eg larger $K$, $\tmax-\tmin$) for weaker or smaller architectures and vice versa. \yx{maybe cut the later part}




\subsection{Experiments on large-scale text-to-image model}
\begin{figure*}[h]
\vspace{-8pt}
    \centering
    \subfigure[FID versus CLIP score]{\includegraphics[width=0.47\textwidth]{img/clip_fid5k_2.pdf}  \label{fig:fid_clip}}\hfill
   \subfigure[FID versus Aesthetic score]{\includegraphics[width=0.47\textwidth]{img/aes_fid5k_2.pdf}    \label{fig:fid_aes}}
     \vspace{-6pt}
    \caption{FID score versus \textbf{(a)} CLIP ViT-g/14 score and \textbf{(b)} Aesthetic score for text-to-image generation at $512 \times 512$ resolution,  using Stable Diffusion v1.5 with a varying classifier-free guidance weight $w=2,3,5,8$. }
      \vspace{-6pt}
\end{figure*}
We further apply Restart to the text-to-image Stable Diffusion v1.5~\footnote{\url{https://huggingface.co/runwayml/stable-diffusion-v1-5}} pre-trained on LAION-5B~\cite{Schuhmann2022LAION5BAO} at a resolution of $512 \times 512$. We employ the commonly used classifier-free guidance~\cite{Ho2022ClassifierFreeDG, Saharia2022PhotorealisticTD} for sampling, wherein each sampling step entails two function evaluations -- the conditional and unconditional predictions. Following \cite{Nichol2021GLIDETP, Saharia2022PhotorealisticTD}, we use the COCO~\cite{Lin2014MicrosoftCC} validation set for evaluation. We assess text-image alignment using the CLIP score~\cite{Hessel2021CLIPScoreAR} with the open-sourced ViT-g/14~\cite{ilharco_gabriel_2021_5143773}, and measure diversity via the FID score. We also evaluate visual quality through the Aesthetic score, as rated by the LAION-Aesthetics Predictor V2~\cite{LAION-AI}. Following \cite{Meng2022OnDO}, we compute all evaluation metrics using 5K captions randomly sampled from the validation set and plot the trade-off curves between CLIP/Aesthetic scores and FID score, with the classifier-free guidance weight $w$ in $\{2,3,5,8\}$. 

We compare with commonly used ODE sampler DDIM~\cite{Song2020DenoisingDI} and the stochastic sampler DDPM~\cite{Ho2020DenoisingDP}. For Restart, we adopt the DDIM solver with 30 steps in the main backward process, and Heun in the Restart backward process, as we empirically find that Heun performs better than DDIM in the Restart. In addition, we select different sets of the hyperparameters for each guidance weight. For instance, when $w=8$, we use $[\tmin,\tmax]{=}[0.1,2], K{=}2$ and $10$ steps in Restart backward process. We defer the detailed Restart configuration to Appendix~\ref{app:config}, and the results of Heun to Appendix~\ref{app:extra-numerical}. 

As illustrated in \Figref{fig:fid_clip} and \Figref{fig:fid_aes}, Restart achieves better FID scores in most cases, given the same CLIP/Aesthetic scores, using only 132 function evaluations (\ie 66 sampling steps). Remarkably, Restart achieves substantially lower FID scores than other samplers when CLIP/Aesthetic scores are high (\ie with larger $w$ values). Conversely, Restart generally obtains a better text-image alignment/visual quality given the same FID.  We also observe that DDPM generally obtains comparable performance with Restart in FID score when CLIP/Aesthetic scores are low, with Restart being more time-efficient. These findings suggest that Restart balances diversity (FID score) against text-image alignment (CLIP score) or visual quality (Aesthetic score) more effectively than previous samplers.

In \Figref{fig:vis}, we visualize the images generated by Restart, DDIM and DDPM with $w=8$. Compared to DDIM, the Restart generates images with superior details~(\eg the rendition of duck legs by DDIM is less accurate) and visual quality. Compared to DDPM, Restart yields more photo-realistic images~(\eg the astronaut). We provide extended of text-to-image generated samples in Appendix~\ref{app:extended-gen}.

\label{exp-t2i}

\begin{figure*}[t]
    \centering
 \subfigure[Restart (Steps=66)]{\includegraphics[width=0.33\textwidth]{img/astronaut_rides_horse_DDIM_restart_True_30.png}}\hfill
    \subfigure[DDIM (Steps=100)]{\includegraphics[width=0.33\textwidth]{img/astronaut_rides_horse_DDIM_restart_False_100.png}}\hfill
   \subfigure[DDPM (Steps=100)]{\includegraphics[width=0.33\textwidth]{img/astronaut_rides_horse_DDPM_restart_False_100.png}}
   \vspace{-6pt}
    \caption{Visualization of generated images with classifier-free guidance weight $w=8$, using four text prompts~(``A photo of an astronaut riding a horse on mars.", "A raccoon playing table tennis",
          "Intricate origami of a fox in a snowy forest" and "A transparent sculpture of a duck made out of glass") and the \textbf{same} random seeds.}
    \label{fig:vis}
    \vspace{-10pt}
\end{figure*}